\chapter{Các giải pháp và đóng góp nổi bật}
\label{ch:contrib}

Trong quá trình thực hiện, một số vấn đề kỹ thuật đáng chú ý đã được giải quyết. Chương này
trình bày các vấn đề đó cùng giải pháp tương ứng.

\section{Thanh toán trực tuyến và tự động hủy đơn quá hạn}

\subsection{Vấn đề}
Khi tích hợp thanh toán trực tuyến, hệ thống cần xử lý trường hợp khách hàng tạo đơn nhưng
không hoàn tất thanh toán. Nếu đơn vẫn giữ tồn kho vô thời hạn, lượng hàng ``ảo'' bị treo sẽ
khiến các khách hàng khác không mua được.

\subsection{Giải pháp}
Hệ thống tích hợp cổng PayOS với quy trình: khi đặt đơn bằng PayOS, một liên kết thanh toán
được tạo kèm thời hạn (mặc định 24 giờ). Kết quả thanh toán được nhận qua webhook và được xác
thực bằng chữ ký HMAC; thao tác đánh dấu ``đã thanh toán'' được thiết kế \textit{idempotent} và
có kiểm tra số tiền để chống xử lý trùng. Một tác vụ cron chạy mỗi 15 phút sẽ quét các đơn PayOS
quá hạn chưa thanh toán, tự động hủy đơn và hoàn trả tồn kho. Nhờ vậy, tồn kho luôn phản ánh
đúng lượng hàng thực sự có thể bán.

\section{Đảm bảo nhất quán dữ liệu khi đặt hàng}

\subsection{Vấn đề}
Nghiệp vụ đặt hàng đụng tới nhiều bảng cùng lúc: kiểm tra và trừ tồn kho của từng biến thể, áp
dụng và tiêu thụ mã giảm giá, ghi nhận đơn hàng. Nếu nhiều khách hàng đặt cùng một sản phẩm
gần như đồng thời, có thể xảy ra tình trạng bán vượt tồn kho (oversell) hoặc dùng voucher vượt
giới hạn.

\subsection{Giải pháp}
Toàn bộ nghiệp vụ đặt hàng được thực hiện bên trong một giao dịch (transaction) của Prisma với
mức cô lập cao. Trong giao dịch, hệ thống kiểm tra lại tồn kho theo thời gian thực, trừ kho, kiểm
tra điều kiện và tiêu thụ voucher một cách nguyên tử. Giá bán được tính lại ở phía máy chủ (ưu
tiên giá flash sale đang hiệu lực) thay vì tin vào dữ liệu gửi từ client, bảo đảm tính toàn vẹn và
chống gian lận giá. Nếu bất kỳ bước nào thất bại, toàn bộ giao dịch được hoàn tác.

\section{Quản lý chương trình khuyến mãi}

\subsection{Vấn đề}
Hệ thống cần hỗ trợ đồng thời hai cơ chế khuyến mãi: mã giảm giá (voucher) áp dụng ở cấp đơn
hàng và flash sale giảm giá ở cấp sản phẩm trong một khoảng thời gian giới hạn.

\subsection{Giải pháp}
Các hàm tính giá thuần (không phụ thuộc trạng thái) được tách riêng và có kiểm thử đơn vị: hàm
tính giá hiệu lực của sản phẩm (ưu tiên giá flash sale đang chạy) và hàm tính mức giảm của
voucher (theo phần trăm hoặc số tiền, có giới hạn mức giảm tối đa). Voucher được kiểm tra điều
kiện (thời hạn, đơn tối thiểu, giới hạn lượt dùng theo tổng và theo người dùng) và tiêu thụ ngay
trong giao dịch đặt hàng, tránh việc một mã bị dùng vượt mức.

\section{Tìm kiếm sản phẩm tiếng Việt không dấu}

\subsection{Vấn đề}
Người dùng Việt Nam thường gõ tìm kiếm không dấu hoặc thiếu dấu (ví dụ ``ao thun'' thay vì ``áo
thun''). Hệ thống cần trả về kết quả phù hợp bất kể người dùng có gõ dấu hay không.

\subsection{Giải pháp}
Hệ thống sử dụng tính năng tìm kiếm toàn văn của PostgreSQL kết hợp tiện ích mở rộng
\texttt{unaccent} để loại bỏ dấu và \texttt{pg\_trgm} để so khớp gần đúng. Nhờ đó, việc chuẩn
hóa dấu được thực hiện thống nhất ở tầng cơ sở dữ liệu, không cần xử lý thủ công ở phía client, cho
kết quả tìm kiếm chính xác và nhất quán.

\section{Bảo mật dữ liệu mức dòng}

\subsection{Vấn đề}
Cần bảo đảm mỗi người dùng chỉ truy cập được dữ liệu của chính mình (đơn hàng, địa chỉ, danh
sách yêu thích), trong khi quản trị viên có quyền truy cập rộng hơn.

\subsection{Giải pháp}
Bên cạnh việc kiểm tra phân quyền ở tầng ứng dụng, hệ thống áp dụng Row-Level Security (RLS)
ở tầng cơ sở dữ liệu PostgreSQL. Các chính sách RLS bảo đảm rằng ngay cả khi có lỗi ở tầng ứng
dụng, dữ liệu của người dùng này cũng không bị lộ cho người dùng khác, tạo thêm một lớp phòng
thủ theo nguyên tắc phòng thủ nhiều lớp (defense in depth).

\section{Quản lý nội dung trang chủ}

\subsection{Vấn đề}
Trang chủ cần hiển thị các banner quảng cáo có thể thay đổi thường xuyên về nội dung và thứ tự,
do quản trị viên tự cập nhật mà không cần lập trình viên can thiệp.

\subsection{Giải pháp}
Hệ thống xây dựng chức năng quản lý banner cho phép thêm, sửa, xóa và đặc biệt là sắp xếp thứ
tự bằng thao tác kéo thả trực quan (sử dụng thư viện dnd-kit). Thứ tự hiển thị được lưu vào cơ sở
dữ liệu và phản ánh ngay trên trang chủ, giúp quản trị viên chủ động trong việc tổ chức nội dung.
